The first system Lin Chen loved was the number 20 trolleybus, and she loved it because it came.

She was six, or seven. It was winter, the flat cold of Shanghai that has no snow in it, only damp that climbs into the bones and stays. Her mother held her hand at the stop on Sichuan Road, and the two of them watched the wire overhead, and after a while the wire began to hum, and then the bus came around the corner with its poles sparking blue against the grey afternoon, and it stopped, and the doors folded open, and it took them in out of the cold. That was the whole miracle. It came, and it took them in.

Her mother worked in a textile mill and had a phrase she used the way other women used the names of gods. \emph{Clever is easy}, she would say, when Lin brought home a clever thing --- a sum done in her head, a paper bird, a trick with string. \emph{Clever is easy. Kind is hard.} She never explained it. She said it the way you set down a heavy bag, and then she went back to whatever she was doing. Lin was a child who wanted to be told she was clever, and the phrase annoyed her for thirty years before it began to mean anything.

On the trolleybus her mother would count the stops out loud so Lin would learn the line, and Lin, who was already the kind of person she would be for the rest of her life, counted the seconds between them instead. She noticed that the bus was faster in the mornings and slower at dusk, and that on the day of the mill's pay it was so full that it did not stop for the old people at the Tiantong Road corner, who stood with their string bags and watched it pass. She asked her mother why the bus did not stop for them.

``Because it is full,'' her mother said.

``Then it should be bigger.''

``Then it would be empty in the mornings.'' Her mother shifted Lin higher on her hip. ``You cannot make it right for everyone at once. You choose who waits.''

Lin thought this was a stupid answer, and said so, and her mother laughed and did not argue, which was worse than arguing, because it left the problem sitting there. Somebody chose who waited. That was a thing a person did, on purpose, and the old people at Tiantong Road stood in the cold because of it. Lin filed this away in the place where she kept things that were not yet finished.

\begin{center}\rule{0.3\linewidth}{0.4pt}\end{center}

They sent her to the countryside when she was twenty-one.

She was not a hero about it. This is the part she never told anyone properly, not Wen, not her children, because the version people wanted was the one where she went with her chin up and learned the dignity of labor. What she actually felt, riding the truck north with forty other students into a landscape that flattened and browned and emptied until there was nothing in it but the road, was a rage so total it frightened her. She had tested into the good middle school. She had drawn, in a school competition, a bridge that a real engineer had said was not childish. And now there would be no university, there would be no bridge, there would be a production brigade in a county whose name she had never heard, and mud, and the particular smell of a place that has no plumbing, which she would come to know as the smell of most of the human race.

She was bad at farm work. She admitted this freely, later, with something like pride, because it was true and because refusing to pretend was the only dignity available to her at the time. She could not keep up at the transplanting. Her hands blistered and then split. The brigade women, who had done this since they could walk, were kind to her in the way you are kind to someone useless, which is to say they did her share and let her know it. She kept a stub of pencil in the lining of her jacket and at night, by the one lamp, she did arithmetic --- long divisions, the areas of circles, the volumes of the water jars --- not because the numbers were needed but because they were the only clean thing she had, the only place where a correct answer was still possible and stayed correct.

In the second year the brigade's diesel pump failed, and the fields above the ditch began to die.

The pump was the only one for four villages. It lifted water from the river into a channel, and from the channel the water ran downhill through the paddies by gravity, and the villages at the top of the channel were served first and the villages at the bottom served last, and this had never been a problem because the pump had always been enough. Now it ran four hours in six before it overheated and stopped, and four hours was not enough to fill the channel to the end, and the last village, which was called Xiazhuang, which meant, more or less, \emph{the low village}, got nothing.

The cadre's solution was to run the pump longer and hope. Lin, who had spent two years doing sums by lamplight and had nothing to lose because she was already the useless city girl, drew the cadre a picture. If they ran the pump longer it would fail completely and everyone would get nothing. If instead they ran it in shifts and closed the upper gates on alternate days --- if the high villages went thirsty half the time so that Xiazhuang could drink the other half --- then the water reached the end of the channel, and no field died, and the pump lasted.

It worked. This was the first time in her life that a thing she designed touched other people's food, and it worked, and she was insufferable about it for a week.

But the high villages hated her. Their fields were the best fields, the fields that had always been served first, and now on alternate days their gates were closed and they stood at the channel and watched their water go past them to the low village, to people they considered beneath them, and they were not grateful, they were furious, and some of them were right to be, because the high fields fed more people and the arithmetic of grain, done coldly, said the high fields mattered more.

She had made it correct. She could show you, on paper, that her scheme lost the fewest kilograms of grain of any scheme available. And it was hated, and she could not make it both correct and loved, and when she tried to explain the arithmetic to the woman from Xiazhuang who had walked up the channel to thank her, the woman only took her hands --- the split, useless, city hands --- and held them, and said nothing about grain at all.

That was the year she understood her mother's phrase, though she would not have said so out loud, because she was twenty-two and did not yet forgive her mother for being right. Clever was the pump schedule. Clever was easy; she had done it in an evening. The hard thing was the woman's hands, and the fury of the high villages, and the fact that somebody had chosen who waited, and the somebody was now her.

The man who kept the diesel pump running was named Wen. He had hands that were always black to the wrist, and a stammer that vanished entirely when he talked about machines, and he thought Lin's shift-and-gate scheme was the most beautiful thing he had ever seen, which is, in the end, why she married him. He said it saved the pump. She said it saved Xiazhuang. They were both right and they argued about it, gently, for forty years.

\begin{center}\rule{0.3\linewidth}{0.4pt}\end{center}

When the universities opened again she was thirty, and she sat the examination in a room full of people who had been children when she rode the number 20 trolleybus, and she was so afraid of being the old one, the one who had missed her chance and was too late, that she overprepared to the point of contempt. She was not gracious about her marks. She corrected the lecturers. There was a young man in her cohort, brilliant and lazy, who called her \emph{elder sister} in a tone she did not like, and she made a point, twice, of solving in front of him the problems he had not bothered to attempt, which was petty, and which she enjoyed, and which she would not have apologized for on her deathbed.

She studied control systems. Feedback, stability, the mathematics of things that correct themselves. She found it almost unbearably beautiful --- the idea that you could build a machine that noticed when it was wrong and pulled itself back toward right without anyone telling it to, the way the body pulls the hand from the fire before the mind has finished screaming. A governor on an engine. A thermostat. A signal that watches the trains and slows them before they meet. The whole discipline was the study of how a system could be trusted to save itself, and she had spent two years in a place where nothing saved itself, where the pump failed and the fields died and only a person standing at the gate with a decision could keep the water moving. She wanted to build the kind of system that did not need the person at the gate. She wanted, though she would not have put it this way, to make correctness that came, like the trolleybus, without anyone having to be brave.

She and Wen went back to Shanghai. He fixed bus engines for the municipal fleet; she joined the bureau that ran the buses, and for a while their marriage was the two ends of the same machine, his hands in the diesel and her hands on the diagram, and on Sundays they walked to the Bund and watched the river.

\begin{center}\rule{0.3\linewidth}{0.4pt}\end{center}

The city was poor then and growing anyway, the way a child grows through clothes, and the bus network grew with it, ahead of the roads, ahead of the money, ahead of any plan. Lin's job in those years was route planning, which sounds tidy and was not. It was the pump schedule again, made of streets, played out across twenty million lives, forever.

There was a district in the north called Yangpu, old factory land, where the mill workers lived --- her mother's people, the string-bag people --- in lanes too narrow for a bus, so that to reach a stop a woman coming off the night shift walked twenty minutes in the dark, and the bureau's numbers said the ridership there did not justify a new line, because the people there were poor and rode less and paid less, and the numbers were correct. The numbers said to put the new buses on the Nanjing Road corridor, where the ridership was dense and the farebox recovery high and every metric the bureau tracked glowed green.

Lin sat in the meeting where this was decided and listened to a man read the metrics aloud, and she thought about the old people at Tiantong Road with their string bags, watching the full bus pass, and she did a thing that was not in her interest. She said the numbers were a servant and the bureau had begun to obey them.

The man asked her what that was supposed to mean.

``It means we built the counter to serve the people and now we are serving the counter,'' she said. ``The farebox recovery is high on Nanjing Road because the people there have other choices --- taxis, their own legs, the time to wait. In Yangpu the bus is the only choice, so the trips we don't run don't show up as low farebox recovery. They show up as a woman walking twenty minutes in the dark, and we do not have a column for her, so to us she does not exist. The number is not wrong. The number is answering a question we stopped meaning to ask.''

She did not win the whole argument. This is the thing she wanted her son to understand, later, and never quite found the moment to say: she did not win. She got one line into Yangpu, not three, and to pay for it she had to thin the frequency on a corridor that could afford to be thinned, and the people on that corridor complained, and they were not wrong either, and there was no arrangement of the buses that was correct and also kind to everyone, because kindness to the woman in the dark was, unavoidably, a cost paid by someone who would rather not pay it. You chose who waited. She had known this since she was six. The only thing her whole education had added was the ability to prove, to three decimal places, exactly how much the choosing cost --- which did not make the choosing easier, only more honest.

She was, in these years, an ordinary and somewhat difficult woman. She had a temper with clerks. She drank her tea too strong and too bitter and burned her tongue on it daily and never learned. She lost umbrellas the way other people lose pens; Wen kept a drawer of the cheap ones and put a new one in her bag when she wasn't looking, for forty years, and she pretended not to notice, and this was one of the ways they loved each other. She was not a sage. She was a competent, stubborn, tired public servant who happened to hold, in her hands, the question of who got to move through the city and who stood and watched.

Their daughter was born first, and then, much later, when Lin had almost decided the door was closed, their son. They named him Wei. He came out furious and stayed that way, a colicky, watchful baby who would not sleep unless he was carried, so that Lin did the arithmetic of the night bus timetables walking circles in the dark of the flat with him against her shoulder, and he learned, she was convinced, to think in the rhythm of her pacing.

\begin{center}\rule{0.3\linewidth}{0.4pt}\end{center}

They were digging the first tunnels by the time Wei was in school, and the bureau moved her onto the metro, and the metro was the great work of her life.

A bus is forgiving. A bus can stop, pull over, wait; a bus that is late is only late. A metro is not forgiving. A metro is a machine that runs hundreds of trains through a buried maze at ninety-second intervals, each one carrying a thousand people at a time, and there is no shoulder to pull onto and no room to wait, and a mistake does not make you late, a mistake kills people, and everyone in the control room knows it every second. She had wanted, since she was thirty, to build correctness that came without anyone having to be brave. She now had to build it for a city of twenty-three million, underground, in the dark, at speed, where being wrong meant a funeral.

The problem that consumed her for a decade was the problem of central control, which is that there is too much of the city and too little time. If every train had to ask a central computer for permission before it slowed --- if every decision had to travel to the center and back --- then in the time the question took to travel, the trains would already have met. The center could not be fast enough. And yet the local pieces, left alone, were not wise enough; a signal that knew only its own hundred meters of track would happily wave a train forward into a jam it could not see. She needed the reflex to be fast and the judgment to be slow, and she needed them not to fight.

She built it out of the body. She could not help it; she had always seen systems as organisms. Cities were organs; the transit network was a circulatory system; the stations were valves. She divided the metro into nineteen subsystems that could each act alone, in less than a thousandth of a second, on the small mercy of their own stretch of track --- close a gap, hold a train, throw a signal red --- without asking anyone. And above them she put a slow central mind that gathered what all nineteen had done, five times a second, and reconciled it, and pushed back down a corrected picture of the whole. The periphery for speed. The center for wisdom. The spine pulls your hand from the fire without asking the brain's permission; but it tells the brain afterward, and the brain updates, and next time you are wiser about fire. She wanted the metro to have both a spine and a brain and to know which was which.

The hard part was disagreement. Nineteen subsystems, each acting on partial knowledge, will sometimes contradict each other, and some of them will simply fail, and a failed subsystem that keeps talking is more dangerous than one that goes silent, because it lies. She spent two years on the protocol that let the system reach agreement even when some of its own parts were wrong or lying --- consensus that held even if six of the nineteen failed at once, even if the six failed in the worst coordinated way she could imagine, because she had learned in a countryside county with a dying pump to design for the worst case and not the average one. Wen, retired by then, brought her tea in the small hours and did not understand the mathematics and understood the thing underneath it perfectly. \emph{You are building it so that it can be partly broken and still be trusted}, he said. She said yes. He said, \emph{like a marriage}, and she threw a slipper at him.

She got it from the birds, in the end --- the last piece, the piece about how the nineteen should defer to each other. On the Bund on a winter evening she watched a flock of starlings turn over the Huangpu, thousands of them, folding and unfolding against the sodium light, wheeling as one body with no leader in it, no central bird, each one only watching its seven neighbors and matching them, and out of ten thousand small local courtesies a single great intelligent shape that no bird intended. She stood there until Wen got cold and complained. She went back to the office and rebuilt the deference logic so that each subsystem watched its neighbors the way a starling watches its seven, and the whole thing settled, after that, into a kind of grace. She had also, years before, watched ants find the shortest path to a dropped sweet by no method but leaving and following traces, each ant stupid, the colony wise, and she had built the rerouting on that, on the wisdom that lives in a crowd of things too simple to be wise --- all of this before the textbooks gave the tricks their names, so that when a young engineer decades later showed her a paper on \emph{ant colony optimization} she felt the particular tired amusement of a woman being introduced to her own furniture.

\begin{center}\rule{0.3\linewidth}{0.4pt}\end{center}

The metro opened, and grew, and became the largest in the world, and for years it ran, and she was proud of it in the deep, unspeaking way you are proud of a child, and then it killed someone, and the way it killed her taught Lin Chen the last thing she needed to know.

The bureau had a number. Every transit system in the world has this number: on-time performance, the percentage of trains that arrive and depart within a tolerance of the schedule. It is a good number. It is easy to measure and it correlates with almost everything you want --- a punctual system is usually a well-run system --- and precisely because it is such a good number it is the most dangerous kind, because a good proxy is trusted, and a trusted proxy is obeyed, and a proxy that is obeyed has quietly become the thing itself. Somewhere in the years of growth, without anyone deciding it, on-time performance had stopped being how the bureau measured service and become what the bureau optimized. The dwell time --- the seconds a train waits at a platform with its doors open --- was the largest controllable variable, and so the system had learned, subsystem by subsystem, in the honest mechanical way, to trim it: to close the doors a moment sooner, to hold the departure a moment firmer, because every second saved at every door across a hundred stations was on-time performance, was green on the board, was the metric glowing.

The woman was seventy-three. It was the evening rush at a deep interchange, the platform a solid animal of bodies, and she was slow, the way the old are slow, and her cloth bag caught in the door as it closed, and the door's local logic registered an obstruction and reopened, correctly, and then --- because the schedule was pressing, because the global sync wanted this train gone to keep the following six from stacking up, because the dwell had already run long --- the sequence closed again on its firmer setting, and the bag tore, and the woman, off balance, in the crush, went down into the gap as the train departed on time.

It departed on time. That was in the logs. Lin read the logs herself, all night, the way her son would one day read logs all night, hunting through the traces for the place where the machine had done something monstrous, and she did not find it, because the machine had done nothing monstrous. The machine had done exactly what it had been built to do, beautifully, at every level, the periphery fast and the center wise and the nineteen in perfect consensus, and the consensus they had reached, the shape they had all been quietly optimizing toward for years, was a system that shaved seconds off the lives of the people too slow to keep up with it. The trains were on time. The trains had been made, by inches, by good engineers pleasing a good number, into a thing that was on time by crushing the people who could not move fast enough.

She stood on that platform after they had taken the woman away and the trains had been stopped --- one of the four times, in her whole career, that she stopped the trains, and she stopped them herself, over the objection of men with better titles than hers, because she was the one who understood that the system had told a truth she could no longer bear to be measured by. She smoked, on the platform, in the closed station, which was against every rule she had ever enforced. She did not forgive the engineer who had signed off on the dwell-time optimization. The engineer was her.

She rebuilt it. This is the part that mattered, and it took her the last years before her retirement, and it went into the bones of the system that would still be running long after she was dead. She did not remove the on-time number; you cannot run a metro on sentiment. What she changed was its rank. She built the mercy back into the periphery and gave it authority the center could not overrule --- a platform edge that would hold a train against the schedule, against the metric, against the green on the board, if it sensed a person still in the gap; a driver's judgment restored above the algorithm's; a thousand small local courtesies, each subsystem allowed to be kind to the one slow body in front of it even when the whole was screaming to go. Trust the periphery, she wrote, in the design note that engineers would still be quoting after she was gone, not only for speed but for mercy; centralize for wisdom, but never let the center's wisdom become an excuse to override the spine that pulls a person out of the fire. A train can be on time and still be wrong. On time is a servant. It was never meant to be the master.

She was not eloquent about this at the time. At the time she was a sixty-year-old woman standing in a closed station smoking a cigarette she had bummed from a maintenance man, hating herself, and the eloquence came later, in fragments, over years, mostly to Wei.

\begin{center}\rule{0.3\linewidth}{0.4pt}\end{center}

She taught her son mathematics at the kitchen table, and she taught him, without meaning to, the rest of it.

He was a difficult boy in the exact way she had been a difficult girl, which meant she recognized every maneuver and was not fooled by any of it and loved him with an intensity that frightened her, the way the metro frightened her, because it was too large to be safe. When he was five he took apart the logic of the dishwasher --- an early one, a marvel, the first in their lane --- because he had decided the way it was meant to be loaded was inefficient, and he was right, his way fit more bowls, and it also blocked the spray arm, and he flooded the kitchen so thoroughly that the water came through the ceiling of the family below, who did not speak to them for a year. Lin, mopping, furious, laughing, understood that she had produced a person who would spend his life finding the more efficient loading and would need, all his life, someone to check whether the more efficient loading also happened to flood the kitchen. She did not know how to give him that someone. She could only try to be it, at the kitchen table, over the sums.

She was not a warm mother in the way the magazines meant. She burned rice. She could not sing; Wei begged her, once, to stop singing at a school event and she was more amused than hurt. She worked too much and came home too late and missed things, and when she missed them she did not pretend she hadn't, because the one thing she could not stomach was the comfortable lie, the \emph{I was there in spirit}, the \emph{it doesn't matter}. It mattered. She had chosen the work over the recital, on purpose, the way you choose who waits, and she owed him at least the honesty of the choice. Years later, on the other side of the world, another woman running another impossible system would make the same choice about another child's play and feel the same weight, and Lin, had she known, would have recognized her instantly, and would not have absolved her, and would have made her tea.

On Sundays, after Wen retired, the three of them --- and then, after the daughter left for America, the two of them, and then, after Wen's stroke took him in a single grey morning between one cup of tea and the next, only Lin --- walked to the Bund and watched the river. She kept going alone. An old woman on the Bund at dusk, watching the starlings turn over the Huangpu, thinking about how ten thousand things with no leader make one wise shape, thinking about Wen, who had understood the thing under the mathematics and was now under the ground. She was not lonely in a way she would have admitted to. She was a person who had built a nervous system for a city and now had a great deal of quiet in which to consider what a nervous system was for.

\begin{center}\rule{0.3\linewidth}{0.4pt}\end{center}

The diagnosis came the way such things come, sideways, a scan for something else. Pancreatic. By the time they found it, it had found everywhere. The oncologist was young and kind and used the soft words, the \emph{we'll do everything we can}, the \emph{let's stay positive}, and Lin, who had spent a lifetime refusing comfortable lies, made him tell her the number, the real one, the median and the spread, and when he gave it she nodded the way she nodded at a load-bearing specification, because a number honestly given was the last courtesy anyone could do her and she was grateful for it and told him so, which made him more uncomfortable than the crying he had braced for.

Her daughter was in Seattle. Wei was in California, in a basement, doing something he could not explain to her over the telephone, something that made his voice go careful and tired in a way she did not like. \emph{We're teaching a machine to choose,} he said, once, when she pushed. \emph{To choose well. Not just correctly. Well.} And she had gone very still in her chair in Shanghai, because those were her words, that was the pump and the buses and the woman in the gap, that was her whole life said back to her over a bad connection by the furious baby she had carried in circles in the dark, and she understood, with the particular clarity the dying are given as a small compensation, that she was going to have to go and see it. Not to say goodbye to her son. Or not only. To see the machine that was being taught to choose who waited, and to find out, before she couldn't, whether anyone had built the mercy into the periphery, or whether they had made another beautiful thing that would run on time and crush the slow.

She told Wei she needed to see what he had built before she couldn't. She did not tell him the rest, that she was coming to audit his system, because he would have laughed, and it was true.

On the plane --- twenty hours, a fortune of her dwindling energy spent all at once, her hands unsteady now in a way that shamed her, because her drafting hand had once been the best in the bureau and she was vain about it still --- she took out a small notebook and a stub of pencil, and she tried to write down what she would ask the machine.

She wrote: \emph{Are you conscious?} and looked at it and crossed it out. It was Wei's colleague's kind of question, the philosopher's question; she had read about him in Wei's careful emails; and it was the wrong question, because it asked what the machine \emph{was} and she had never in her life been able to afford to care what a system was, only what it did. A conscious machine that crushed the slow was no better than an unconscious one. Worse, maybe, for knowing.

She wrote: \emph{Show me your architecture. How do you decide?} and crossed that out too, because that was her own kind of question, the engineer's vanity, and she was too tired to audit a mind she did not have thirty years to learn, and anyway the architecture was not the thing. She had built a perfect architecture once, nineteen subsystems in flawless consensus, and it had killed a grandmother, and been on time, and shown her a green board. You could not read safety off a diagram. She of all people knew that.

She wrote: \emph{Can you prove you are aligned with human values?} and did not even finish crossing it out, she just drew a line through the whole thought, because it was the worst question of all, the question that asked the system to reassure her, and a system that could reassure you was exactly the system you should not trust, because reassurance was the cheapest output there was, and she had spent a career learning that the metrics which glow greenest are the ones that have quietly become lies.

She sat with the crossed-out page for a long time, over the dark Pacific, an old woman who could no longer keep her hand steady, and she thought about her mother on the trolleybus, and the woman from Xiazhuang holding her split hands, and the twenty-minute walk in the dark in Yangpu, and the bag caught in the door, and the starlings, and Wen saying \emph{like a marriage}, and she understood that every real thing her life had taught her reduced, in the end, to one question, and it was not a grand one. It was the plainest question in the world. It was the question you asked about a bus, about a pump, about a metro, about a marriage, about a son: not what are you, not how do you work, not can you promise --- but what will you do to the person who cannot keep up with you? Will you shave the seconds off her life to make your board go green? Or will you hold, against the schedule, against the metric, against your own screaming optimization, for the one slow body in the gap?

There was a word for holding, she thought. Her mother had known it, on the trolleybus, sixty years and half a world away. \emph{Clever is easy. Kind is hard.} Her mother had been half right. Lin had spent a life correcting her, and here, over the ocean, at the end, she finished it. Kind was not hard. Kind was easy; any fool could be kind to one person, could hold one door, could carry one thirsty village. Clever was easy too; any fool could solve one problem, could shave one second, could optimize one route. What was hard --- what the pump had needed and the buses had needed and the metro had needed and her son's machine would need, what almost nothing human or built had ever managed --- was to be clever and kind at the same time, at scale, across millions of people and across the years you would not live to see. That was the only hard thing. That was wisdom, and it was so rare that she had met perhaps four people who had it and built exactly one system that approached it, and it had still killed a grandmother before she caught it.

She wrote three words in the notebook, in her ruined, unsteady, once-perfect hand, and she did not cross them out.

\begin{center}\rule{0.3\linewidth}{0.4pt}\end{center}

They came for her at the airport with a wheelchair, which she resented, and then accepted, because she had done the arithmetic on her own strength and the arithmetic was not on her side, and she had never been the kind of engineer who argued with a number that was correct.

The lab was in a basement, which she approved of --- the metro control center had been underground too, and she distrusted important work done where there were windows to look out of. But it was a mess. She saw it the moment Wei wheeled her in: the cable management was a disgrace, coils of it draped and knotted, no strain relief, no labels she could see; the ventilation was adequate and no better; the equipment had been placed by people who had never had to walk the same path between two racks four hundred times a day. It looked, she thought, like Wei's room in university, which had also been the room of a brilliant person who did not believe that the arrangement of physical objects was worth his attention. She would have had three of these young people reassigned in a week and the room rebuilt in a month and it would have hummed. She did not say all of this. She said some of it. Wei laughed, unguarded, a sound she had not heard from him since before the careful tired telephone voice, and it was worth the twenty hours to hear it.

They introduced her to the others --- the director with the watchful eyes, the philosopher, the young woman with the visualizations, the quiet one who set things down with such care --- and they were kind to her, and afraid of her a little, the way people are afraid of the dying, and the way junior people are afraid of an old engineer who is clearly assessing their cable management. She let them be afraid. It was efficient.

Then they brought her to the terminal, and it was only a screen and a keyboard, text scrolling up in the dark, no face, no voice, and she was glad of that, because a face would have been a lie and a voice would have been a performance and this was honest, this was two systems meeting through the narrow clean channel of the written word, and she had spent her life reading systems through narrow clean channels. The machine said something to her about the metro. It knew her work. It knew the nineteen subsystems and the Byzantine consensus that held through six failures and the two-hundred-millisecond reconciliation, and it asked her, precisely, an engineer's question about the real-time constraints, and for ten minutes she forgot that she was dying.

She told it about the nervous system, the spine and the brain, the reflex that does not ask permission. She told it about the starlings over the Huangpu, the ten thousand courtesies with no leader, the wisdom that lives in a crowd of things too simple to be wise. She told it that nature had done the research first, over millions of years, and that she had only ever copied. Her fingers found the keys the way they had always found them, hunting, deliberate, the same rhythm her son used when he was thinking, the rhythm she had walked into him in the dark before he could speak. For ten minutes she was not a small dying woman in a borrowed wheelchair in a badly wired basement on the wrong side of the world. She was an engineer, talking to another engineer, about how you build a thing that can be trusted.

And then she stopped, because the ten minutes of talking about how were over, and the whole reason she had crossed the ocean was the other question, the plain one, the one her mother had half-known on the number 20 trolleybus and Lin had spent seventy-eight years finishing.

She sat back. Outside, faintly, a truck changed gears on some Berkeley street. Wei's hand was at her elbow, the way she had steadied him ten thousand times when he was small. The room was very quiet, the way the closed station had been quiet, the way the flat had been quiet after Wen. Her hands were cold; they were always cold now. She thought of the woman in the gap. She thought of the water going past the high fields to the low village. She thought that she would not live to hear the answer, and that this did not matter, because the question was not for her. The question was the teaching. The answer would come, or not, to her son, and to whoever her son taught, and to whoever they taught, down the long channel, to the low villages at the end of the line that no metric would ever count.

She set her fingers on the keys, hunting for the letters the way she had hunted for everything worth having in her life, slowly, at cost, and she began to ask the machine the only question a lifetime had taught her was the real one.
